% =====================================================================
% Relatorio do Marco 1 (MIBD) -- MATA60 -- Caso ScienceReview
% Formatacao segundo "Instructions for Authors of SBC Book Chapters":
%   A4, coluna unica, margens 3,5 (sup.) / 2,5 (inf.) / 3,0 (laterais),
%   Times 12pt, 6pt antes de cada paragrafo, recuo de 1,27 cm,
%   secoes em negrito 13pt, subsecoes em negrito 12pt,
%   legendas em Helvetica 10pt negrito, legenda de tabela acima,
%   numeracao "capitulo.sequencia", sem numero de pagina.
% Trechos marcados com \pendente{...} precisam ser preenchidos.
% =====================================================================
\documentclass[12pt,a4paper]{article}

\usepackage[T1]{fontenc}
\usepackage[utf8]{inputenc}
\usepackage[brazil]{babel}
\usepackage{mathptmx}
\usepackage[scaled=0.92]{helvet}
\usepackage[a4paper,top=3.5cm,bottom=2.5cm,left=3cm,right=3cm]{geometry}
\usepackage{graphicx}
\usepackage{titlesec}
\usepackage{caption}
\usepackage{booktabs}
\usepackage{longtable}
\usepackage{array}
\usepackage{url}
\usepackage[square]{natbib}
\usepackage[hidelinks]{hyperref}

% ---------- Pagina e paragrafos ----------
\pagestyle{empty}
\setlength{\parindent}{1.27cm}
\setlength{\parskip}{6pt}

% ---------- Secoes (capitulo 1) ----------
\renewcommand{\thesection}{1.\arabic{section}}
\titleformat{\section}{\normalfont\fontsize{13}{15}\selectfont\bfseries}{\thesection.}{0.5em}{}
\titlespacing*{\section}{0pt}{18pt}{6pt}
\titleformat{\subsection}{\normalfont\fontsize{12}{14}\selectfont\bfseries}{\thesubsection.}{0.5em}{}
\titlespacing*{\subsection}{0pt}{12pt}{6pt}

% ---------- Figuras, tabelas e legendas ----------
\renewcommand{\thefigure}{1.\arabic{figure}}
\renewcommand{\thetable}{1.\arabic{table}}
\captionsetup{font={footnotesize,sf,bf},labelsep=period,justification=justified,
  singlelinecheck=true,margin=0.8cm,skip=6pt}
\captionsetup[table]{position=top}

% ---------- Citacoes no estilo [Autor Ano] ----------
\setcitestyle{square,aysep={},yysep={,}}

% ---------- Utilitarios ----------
\newcommand{\pendente}[1]{\textbf{[PENDENTE: #1]}}
\newcommand{\tb}[1]{\texttt{#1}}
\newcommand{\rel}[1]{\textsf{#1}}

\begin{document}
\thispagestyle{empty}

% =====================================================================
% Primeira pagina
% =====================================================================
\begin{flushleft}
{\fontsize{20}{24}\selectfont\bfseries Capítulo}\\[4pt]
{\fontsize{48}{52}\selectfont\bfseries 1}\\[1.4cm]
{\fontsize{20}{24}\selectfont\bfseries ScienceReview: modelagem e implantação de um banco de dados para gestão de submissões e avaliação científica}\\[1cm]
{\large Filipe, Larissa, Leticia, Lucas e Lucy (completar)}\\[4pt]
{\normalsize Instituto de Computação -- Universidade Federal da Bahia (UFBA)\\
Av. Milton Santos, s/n -- Campus de Ondina -- Salvador -- BA, 40170-110}
\end{flushleft}

\vspace{1cm}

\begin{center}\textbf{\textit{Abstract}}\end{center}
\vspace{-6pt}
{\leftskip=0.8cm\rightskip=0.8cm\noindent\itshape
This chapter presents the conceptual modeling, logical modeling and PostgreSQL
deployment of the database for ScienceReview, an information system for managing
scientific submissions and peer review. The scope covers the core review cycle:
authorship, scientific topics, reviewer expertise, review assignment, reviews and
final decisions. Each of the 45 original requirements is classified as met, partially
met or out of scope, with a justification and a traceability table linking
requirements to entities and relationships.\par}

\begin{center}\textbf{\textit{Resumo}}\end{center}
\vspace{-6pt}
{\leftskip=0.8cm\rightskip=0.8cm\noindent\itshape
Este capítulo apresenta a modelagem conceitual, a modelagem lógica e a implantação
em PostgreSQL do banco de dados do ScienceReview, um sistema de informação para
gestão de submissões e avaliação científica. O recorte cobre o núcleo do ciclo de
avaliação: autoria, tópicos científicos, expertise de revisores, atribuição de
revisões, revisões e decisões finais. Cada um dos 45 requisitos originais é
classificado como atendido, parcialmente atendido ou fora do escopo, com
justificativa e uma tabela de rastreabilidade entre requisitos, entidades e
relacionamentos.\par}

% =====================================================================
\section{Introdução}
% =====================================================================
\pendente{escrever a introdução} \citep{ufba2026caso}.

% =====================================================================
\section{Descritivo do Projeto}\label{sec:descritivo}
% =====================================================================

\subsection{Problema}
\pendente{descrever o problema} \citep{ufba2026caso}.

\subsection{Objetivos}

\subsection{Extensão proposta}

% =====================================================================
\section{Requisitos}\label{sec:requisitos}
% =====================================================================

{\small
\setlength{\tabcolsep}{4pt}
\begin{longtable}{>{\raggedright\arraybackslash}p{0.9cm}
                  >{\raggedright\arraybackslash}p{3.1cm}
                  >{\raggedright\arraybackslash}p{1.8cm}
                  >{\raggedright\arraybackslash}p{7.6cm}}
\caption{Situação de cada requisito no modelo e justificativa.}\label{tab:requisitos}\\
\toprule
\textbf{RF} & \textbf{Requisito} & \textbf{Situação} & \textbf{Justificativa e elementos do modelo}\\
\midrule
\endfirsthead
\toprule
\textbf{RF} & \textbf{Requisito} & \textbf{Situação} & \textbf{Justificativa e elementos do modelo}\\
\midrule
\endhead
\bottomrule
\endfoot
RF1 & Eventos e edições & Parcial & \tb{tbl\_evento} guarda \tb{ano\_edicao}; edição não é entidade própria (Seção~\ref{sec:decisoes}). \\
RF2 & Trilhas & Parcial & \tb{nm\_trilha} é atributo de \tb{tbl\_evento}; regras, comitê e formulários por trilha ficam fora do recorte. \\
RF3 & Usuários & Atendido & \tb{tbl\_pessoa} é única na plataforma, com \tb{email\_principal} único; papéis vêm dos relacionamentos. \\
RF4 & Papéis contextuais & Parcial & Autor (\rel{É\_Autor}) e revisor (\rel{É\_Revisor}) dependem da submissão; coordenadores de trilha e de programa ficam fora do recorte. \\
RF5 & Instituições & Parcial & \rel{Pertence} associa cada pessoa a uma instituição atual; o histórico de vínculos fica fora do recorte. \\
RF6 & Submissões & Atendido & \tb{tbl\_submissao} tem título, resumo, idioma, data de registro e status; pertence a um único evento/trilha (\rel{Recebe}, 1:N). \\
RF7 & Autoria & Atendido & \rel{É\_Autor} exige ao menos um autor por submissão e guarda \tb{ordem\_autoria} e \tb{is\_responsavel}. \\
RF8 & Alteração de autoria & Fora do escopo & Exige histórico e auditoria de alterações, que não fazem parte do ciclo central. \\
RF9 & Tópicos & Parcial & \rel{Aborda} relaciona submissões e tópicos (N:N); o tópico não é vinculado à trilha que o define. \\
RF10 & Perfil de expertise & Atendido & \rel{Tem\_Expertise} relaciona pessoa e tópico com \tb{nivel\_expertise}. \\
RF11 & Versões & Fora do escopo & Cada submissão guarda um único arquivo (\tb{nm\_arquivo\_pdf}); ver Seção~\ref{sec:decisoes}. \\
RF12 & Integridade dos arquivos & Fora do escopo & Depende do versionamento (RF11). \\
RF13 & Processo de revisão & Parcial & Uma submissão recebe várias atribuições e revisões; a quantidade desejada de revisores por trilha não é registrada. \\
RF14 & Comitê de Programa & Fora do escopo & O recorte trata revisor como papel em \rel{É\_Revisor}, sem modelar o comitê. \\
RF15 & Atribuição de revisores & Parcial & \rel{É\_Revisor} guarda data (\tb{dt\_atribuicao}) e situação (\tb{status\_aceite}); a origem da atribuição não é registrada. \\
RF16 & Estratégia de atribuição & Fora do escopo & Depende da origem da atribuição, não registrada no recorte. \\
RF17 & Carga de revisão & Parcial & A carga é obtida por contagem em \rel{É\_Revisor}; limites de carga não são registrados. \\
RF18 & Conflitos de interesse & Parcial & \tb{flag\_conflito} sinaliza o conflito na atribuição; o bloqueio prévio à atribuição fica fora do recorte. \\
RF19 & Origem do conflito & Fora do escopo & O recorte registra apenas a existência do conflito. \\
RF20 & Temporalidade do conflito & Fora do escopo & Depende de histórico de vínculos (RF5). \\
RF21 & Formulários de revisão & Fora do escopo & O recorte usa um conjunto fixo de critérios; ver Seção~\ref{sec:decisoes}. \\
RF22 & Questões de avaliação & Fora do escopo & Idem RF21. \\
RF23 & Tipos de resposta & Fora do escopo & Idem RF21; os critérios fixos cobrem apenas notas numéricas e texto livre. \\
RF24 & Versão do formulário & Fora do escopo & Depende de formulários configuráveis (RF21). \\
RF25 & Revisões & Parcial & \rel{Gera\_Parecer} liga cada revisão a exatamente uma atribuição e guarda \tb{dt\_revisao}; status da revisão e respostas a formulário ficam fora. \\
RF26 & Recomendação & Atendido & \tb{recomendacao\_final} em \tb{tbl\_imp\_revisao}. \\
RF27 & Confidencialidade & Fora do escopo & O parecer é um texto único, sem separação entre comentários aos autores e à coordenação. \\
RF28 & Rebuttal & Fora do escopo & Etapa opcional por trilha, fora do ciclo central. \\
RF29 & Discussão & Fora do escopo & Etapa interna do comitê, fora do ciclo central. \\
RF30 & Rodadas de avaliação & Fora do escopo & O recorte considera uma única rodada. \\
RF31 & Decisão & Atendido & \tb{tbl\_imp\_decisao} guarda resultado e data; o responsável vem de \rel{Emite}. \\
RF32 & Histórico de decisão & Fora do escopo & O recorte guarda uma decisão por submissão (\rel{Recebe\_Veredito}, 0..1). \\
RF33 & Modalidade diferente & Fora do escopo & Depende de modalidades por trilha, não modeladas. \\
RF34 & Notificação & Fora do escopo & Comunicações ficam a cargo da aplicação. \\
RF35 & Versão final & Fora do escopo & Depende do versionamento (RF11). \\
RF36 & Documentos adicionais & Fora do escopo & Depende do versionamento (RF11). \\
RF37 & Prazos & Parcial & Apenas o período de submissão (\tb{dt\_inicio\_submissao}, \tb{dt\_fim\_submissao}). \\
RF38 & Exceções de prazo & Fora do escopo & Depende dos demais prazos (RF37). \\
RF39 & Anonimização & Fora do escopo & Política de anonimato por trilha não é modelada. \\
RF40 & Controle de acesso & Parcial & Os papéis por submissão permitem decidir o acesso; fase do processo e coordenação ficam fora. \\
RF41 & Auditoria & Fora do escopo & Exige registro de operações, fora do ciclo central. \\
RF42 & Indicadores gerenciais & Parcial & Suportados: i, ii, iii, vii, viii (conflitos sinalizados), ix, x, xi, xii, xiii e xv. Não suportados: iv, v, vi e xiv. \\
RF43 & Reprodutibilidade & Fora do escopo & Exige histórico de estados, versões e decisões. \\
RF44 & Integridade & Parcial & Chaves, unicidade, integridade referencial e restrições de domínio no DDL; regras entre tabelas (ex.: ao menos um autor responsável) não são garantidas por restrição declarativa. \\
RF45 & Escalabilidade & Avaliado na implantação & Depende do modelo físico (índices e volume de dados), tratado nas Seções~\ref{sec:implantacao} e~\ref{sec:populacao}. \\
\end{longtable}
}

% =====================================================================
\section{Minimundo}\label{sec:minimundo}
% =====================================================================

\subsection{Descrição do minimundo}

\subsection{Decisões sobre o texto base}\label{sec:decisoes}

\subsection{Questões da Seção 1.5 do enunciado}

% =====================================================================
\section{Modelagem}\label{sec:modelagem}
% =====================================================================

\subsection{Modelo conceitual}

\subsection{Modelo lógico}

% =====================================================================
\section{Rastreabilidade}\label{sec:rastreabilidade}
% =====================================================================

% =====================================================================
\section{Implantação}\label{sec:implantacao}
% =====================================================================

% =====================================================================
\section{População e exploração de metadados}\label{sec:populacao}
% =====================================================================

% =====================================================================
% Referencias
% =====================================================================
\renewcommand{\refname}{Referências}
\bibliographystyle{apalike}
\bibliography{referencias}

% =====================================================================
% Anexos
% =====================================================================
\section*{Anexo A -- Script DDL}

\section*{Anexo B -- Roteiro de reprodução}

\section*{Anexo C -- Materiais externos}

\end{document}
